% =====================================================================
%  GMN3 §11 量子和乐 + §12 热带标签（含 §12.1 $N=3$ Argyres--Douglas 例）
%  原文：Lamination-Paper.tex 5864--6289 行；PDF p.99--106（印刷 p.98--105）
%  编号公式 (11.1)--(11.3)、(12.1)--(12.12) 共 15 条；
%  插图 4 幅（cylinder / Flip-Traffic / tropic-onto-A / tropic-onto-B）
%  跨块引用（节号、式号、附录）暂写原文编号并加 % TODO，合稿时统一改 \ref/\eqref
% =====================================================================

\section{量子和乐}\label{sec:Quantum-Holonomy}

在第 10 节% TODO 合稿改 \ref{sec:Examples}
中，我们说明了如何计算 $A_1$ 理论的带框 BPS 简并度。本节我们勾勒如何——原则上——更进一步，算出 $A_1$ 理论完整的保护自旋特征标，从而算出这些例子中的形变函数代数。关键想法是把穿墙公式与第 10 节% TODO 合稿改 \ref{sec:Examples}
开头所解释的提升原理结合起来。

\insfig[0.55\textwidth]{cylinder}{在 $(u,\zeta)$（以及各耦合常数）的参数空间中，凡 $T_\WKB(u,\zeta)$ 形如本图之处，我们都可以用提升原理来计算与曲线 $\wp$ 相伴的线算子的 PSC。}

在一些特殊情形下，提升原理可以直接了当地应用（第 10 节% TODO 合稿改 \ref{sec:Examples}
中已经顺带这样做过）。设有一条曲线 $\wp$ 与一个三角剖分，使得 $\wp$ 的局部邻域形如图 \ref{fig:cylinder}。设我们想计算线算子 $L_\zeta(\wp)$ 的 PSC 在 $y=1$ 处的特殊化（这里表示取 $\CR=\mathbf{2}$ 是默认的）。也就是说，我们想计算期望值 $\langle L_\zeta(\wp)\rangle'$。应用附录 F % TODO 合稿改 \ref{app:traffic}
的 traffic rule，得到
\begin{equation}\label{eq:Prime-Wilson-Spec}
\langle L_\zeta(\wp)\rangle' = \sqrt{\tCY_+\tCY_-} + \frac{1}{\sqrt{\tCY_+\tCY_-}} + \sqrt{\frac{\tilde\CY_+}{\tCY_-}}.
\end{equation}
显然可以把提升原理应用于这个表达式，于是断定
\begin{equation}\label{eq:PromPrinc-1}
F(L_\zeta(\wp)) = X_{\half(\gamma_+ + \gamma_-)} + X_{-\half(\gamma_+ + \gamma_-)} + X_{\half(\gamma_+ - \gamma_-)}.
\end{equation}

现在设我们想对某个一般的 $(u,\zeta)$ 计算 $F(L_\zeta(\wp))$。数据 $(u,\zeta)$（以及各耦合常数）决定了一个 Seiberg--Witten 微分，以及与之相伴的 WKB 三角剖分 $T_\WKB(\vartheta,u)$，其中 $\vartheta=\arg\zeta$，详见文献 \cite{Gaiotto:2009hg} 第 6 节。通过变动 $u,\zeta$ 与理论的各耦合常数，我们试图找到从该 WKB 三角剖分到图 \ref{fig:cylinder} 那种形状的一个态射。若能找到，则对连接这两个三角剖分的一连串 flip 依次应用变换规则 (3.47)% TODO 合稿改 \eqref{eq:Spin-KS-tmn-S}
、(3.48)% TODO 合稿改 \eqref{eq:Spin-KS-ii}
，便可从 \eqref{eq:PromPrinc-1}
出发，得到 $(u,\zeta)$ 处的保护自旋特征标。

上述计算 PSC 的算法有两个重要的漏洞。第一，只有当 $C$ 容许图 \ref{fig:cylinder} 那种形状的三角剖分时，该算法才管用。现实中确实存在没有这种三角剖分的情形。所幸我们可以借助 J. Teschner 在 Teichm\"uller 空间量子化中所用的一个步骤来处理这种情形 \cite{Teschner:2005bz}。简述如下：取 $C$ 的一个 trinion（三孔球面）分解，使得 $\wp$ 是该分解中的一条切割闭链；各 trinion 的边界被标为「ingoing」与「outgoing」；存在一个与该 trinion 分解相容的三角剖分，用来在 Teichm\"uller 空间上定义 Fock 坐标（见文献 \cite{Teschner:2005bz} 第 13、14 节）。我们上面描述的算法适用于文献 \cite{Teschner:2005bz} 式 (15.4) 的情形；另一种情形出现在该文式 (15.5)。不过，正如该文所述，可以调用一个递推步骤来计算 $\langle L_\zeta(\wp)\rangle'$，而提升原理将适用于由此递推步骤所得的结果。

第二个漏洞是：即使存在图 \ref{fig:cylinder} 那种形状的三角剖分，也并不能一眼看出存在这种形状的 WKB 三角剖分。我们不认为这是一个严重的漏洞，因为总可以选取适当的弱耦合区域，使 $\wp$ 的管状邻域是一根细长的管子，而 $L_\zeta(\wp)$ 对应于规范群中相应 $SU(2)$ 因子的基本 Wilson 圈。在这样的参数空间区域里，WKB 三角剖分将会是合适的那一型，在我们看来相当合理；但本文不给出这一断言的详细论证。

我们还要再作几点评注：

\begin{enumerate}
\item 我们在第 4.4 节% TODO 合稿改 \ref{subsec:formal-SU2}
中已经看到，在 $SU(2)$、$N_f=0$ 的情形，一些强正的形式线算子可以纯粹代数地轻易算出；而第 5 节% TODO 合稿改 \ref{sec:Cluster-Algebras}
所述的与簇代数的联系表明，可以用局部规则来拼接模空间中的各个腔。于是，利用提升原理就能相当高效地算出某些 PSC。

\item 上述计算 PSC 的规则在某种意义上是 $C$ 上的非局部规则。我们希望能有计算带框保护自旋特征标的\emph{局部}规则，类似于我们在 $y=+1$ 处拥有的 traffic rule。我们预期这类局部规则将来自一个精细的 2d-4d 穿墙公式，其存在性曾在文献 \cite{Gaiotto:2009fs} 中被提示。这项工作正在进行中。\footnote{A. Goncharov 告诉我们，在簇系综（cluster ensembles）理论中，已有一些已知但尚未发表的规则可计算类似量 \cite{MR2567745}。}

\item 由与 Teschner 工作的这一联系可知：$A_1$ 理论中 $\CM$ 上的形变函数代数，自然地同构于 $C$ 的量子 Teichm\"uller 理论中量子测地算子之代数（后者量子化的是常曲率 $-1$ 度量下的测地长度）。我们要强调，我们的框架纳入了不规则奇点与叶层，因而为 Teichm\"uller 语境中已有的工作提供了有意思的推广。

\item 也许还可以更进一步，说明形式展开 (3.38) % TODO 合稿改 \eqref{eq:form-gen}
可以借助对 $\CM$ 的某个适当实切片作某种量子化而获得具体含义，即
\begin{equation}\label{eq:q-proposal}
\hat L_\wp = \sum_\gamma \fro(u,L,\zeta,\gamma;y) X_\gamma
\end{equation}
其中 $X_\gamma$ 现在是作用在某个确定的希尔伯特空间上的具体量子算子。事实上，文献 \cite{Teschner:2005bz} 以及 \cite{MR2567745,clqd2,qd-cluster} 都构造出了 $X_\gamma$ 的具体算子解释。不过，\eqref{eq:q-proposal} 应当有更广的适用范围，且不应依赖于实切片的选择。在选取某些实切片与特定量子化方案之后，我们或许可以期待与 Nekrasov--Shatashvili \cite{Nekrasov:2009ui,Nekrasov:2009rc} 以及 Nekrasov--Witten \cite{Nekrasov:2010ka} 的工作建立某种联系。
\end{enumerate}

\section{热带标签}\label{sec:Tropic}

本节我们勾勒一种用红外数据标记 $\CN=2$ 理论中单线算子的办法。

寻找这样的标记有三重动机。第一，这种标记是可取的，因为迄今我们只能在具有拉氏量的理论中给出标记（见第 2.3 节% TODO 合稿改 \ref{subsec:LabelsForLines}
）。红外标记同样可以用于非拉氏理论。第二，刻画 $\langle L\rangle$ 在大 $R$ 或大 $\zeta$ 时的渐近行为是有意义的。这有数学上的动机，因为我们知道和乐的迹是 Hitchin 模空间上有意思的函数；因此考察 $\zeta\to 0,\infty$ 的渐近行为很有意思，而结果表明它非常微妙。这里似乎还有一些有意思的物理学应用。第三，我们想补上第 6 节% TODO 合稿改 \ref{sec:linehol}
中的一个缺口：我们想证明存在足够多的不同线算子 $L_i$，足以「反解」展开 (6.3)% TODO 合稿改 \eqref{eq:trh-large-r}
，从而把 $\CY_\gamma$ 用诸 $L_i$ 表示出来。

结果表明，「热带代数簇」这一课题对处理这些问题很有用。我们将用到 Fock 与 Goncharov 的工作 \cite{MR2233852} 中迄今尚未在我们的故事里扮演过角色的某个方面。

关键想法是考察某个单线算子的 vev $\langle L\rangle$ 的首项渐近行为，例如在 $\zeta\to 0$ 时。达布展开 (6.3)% TODO 合稿改 \eqref{eq:trh-large-r}
把 $\langle L\rangle$ 表示成诸 $\CY_\gamma$ 之和。当 $\zeta$ 沿某条射线趋于 $0$，或固定 $(u,\zeta)$ 而令 $R\to\infty$ 时，和式中会有一项占主导。若考虑 $\widehat\CB\times\widehat\IC^*$ 中的一条路径 $(u_s,\zeta_s)$，则两项 $\fro(L_1,\gamma_1)\CY_{\gamma_1}$ 与 $\fro(L_2,\gamma_2)\CY_{\gamma_2}$ 只可能在路径上满足 ${\Re}(Z_{\gamma_1}/\zeta)={\Re}(Z_{\gamma_2}/\zeta)$ 的点处交换主导地位。对单线算子而言，从 halo 图像出发自然预期：这样两项必定满足 $\gamma_1-\gamma_2=\gamma$，其中 $\gamma$ 支撑一个「原味」BPS 态。因此，对 $\gamma\in\Gamma$，我们定义\emph{反 BPS 墙}（anti-BPS wall）$\check W(\gamma)$ 为这样一道墙：\footnote{在我们唯一的例子即 $A_1$ 理论中，这一定义将被证明是正确的。但我们注意到：尽管它是自然的，却并不能从 halo 图像严格推出「发生主导地位交换的墙只有这些墙」。原则上可能出现 $\gamma_1-\gamma_2$ 不支撑「原味」BPS 态的情形。若真有这类例子，下文关于热带理论的定义就得作修改。}
\begin{equation}
\check W(\gamma):= \{(u,\zeta): Z_\gamma/\zeta \in -i\IR_+ \qquad {\rm and} \qquad \Omega(\gamma;u)\not=0\}.
\end{equation}

我们之所以选取 $-i$ 这个符号，下文自会明了。这些墙的补集有若干连通分支，它们被定义为\emph{反腔}（anti-chamber），记作 $\check c$。一般地我们预期：在给定的反腔内，达布展开 (6.3)% TODO 合稿改 \eqref{eq:trh-large-r}
中会有单独一项主导渐近行为，于是可以写成
\begin{equation}
\langle L\rangle \sim \fro(L;\gamma_t)\exp(\pi R Z_{\gamma_t}/\zeta)\left(1+\cdots\right)
\end{equation}
其中 $\gamma_t\in\Gamma_\CL$ 是某个向量。我们把它称为\emph{热带标签}。一般地它依赖于 $L$ 与 $\check c$，故记作 $\gamma_t(L,\check c)$。注意我们有「热带公式」：
\begin{equation}\label{eq:tropic-mult-line}
\gamma_t(LL',\check c) = \gamma_t(L,\check c) + \gamma_t(L',\check c)
\end{equation}
以及
\begin{equation}
{\rm Re}(Z_{\gamma_t(L+L',\check c)}/\zeta) = {\rm Max}[{\rm Re}(Z_{\gamma_t(L,\check c)}/\zeta), {\rm Re}(Z_{\gamma_t(L',\check c)}/\zeta)].
\end{equation}
（关于 \eqref{eq:tropic-mult-line}：乘积 $LL'$ 一般不是单线算子，但可以按单线算子分解；所得诸算子中有一个会具有 \eqref{eq:tropic-mult-line} 给出的热带标签。）

更形式地说，我们把一个 $\CN=2$ 理论称为\emph{热带的}，若存在 $\Gamma$ 的一个（依赖于反腔的）基 $\{\gamma_a\}$，使得\footnote{我们也可以改用 $\langle L_\zeta\rangle'$ 与 $\widetilde\CY_\gamma$ 来表述，其优点是一切展开系数都为正，因而不可能发生相消。}

\begin{enumerate}
\item 在任一反腔 $\check c$ 中，对一切 $a$，当 $R\to\infty$ 与 $\zeta\to 0$ 时的渐近行为都满足 $\CY_a\to\infty$。此外，一切单线算子都具有形式：
\begin{equation}
\langle L\rangle = \fro(L,\gamma_t)\CY_{\gamma_t}P(1/\CY_a)
\end{equation}
其中 $P$ 是一个最低阶项为 $1$ 的多项式，而 $\gamma_t\in\Gamma_\CL$ 在反腔 $\check c$ 内为常数。（多项式 $P$ 可以在反腔内改变。例如，通常的腔会把该反腔分成不止一个连通分支。）

\item 跨过反 BPS 墙 $\check W(\gamma_0)$ 时，我们有变换律
\begin{equation}\label{eq:Tropical-Transformation}
\gamma_t(L_i,\check c') = \gamma_t(L_i,\check c) - \Omega(\gamma_0;c)\langle\gamma_t(L_i,\check c),\gamma_0\rangle_+ \gamma_0
\end{equation}
（该反墙位于腔 $c$ 内部）。

\item 对一切 $\gamma\in\Gamma_\CL$ 与任一反腔 $\check c$，都存在一个单线算子 $L_\gamma$，使得 $\gamma=\gamma_t(L_\gamma,\check c)$。

\item 若 $\gamma_t(L,\check c)=\gamma_t(L',\check c)$，则 $L=L'$。
\end{enumerate}

用物理的话说，这表示：一个单线算子通常被赋予一个紫外标签（例如 $\CL$ 中的一个向量），而它也可以由该理论（扩张的）荷格 $\Gamma_\CL$ 中的一个红外标签来唯一标记。不过，与紫外标签不同，红外标签依赖于红外参数（通过其对 $\check c$ 的依赖），并且会经历穿墙。

Fock 与 Goncharov 的工作 \cite{MR2233852} 提示如下猜想：$\CS$ 类中的所有理论都是热带的。下面我们勾勒为什么 $A_1$ 理论确实是热带的。（以下论证严重依赖文献 \cite{Gaiotto:2009hg} 所发展的技术，我们假定读者熟悉该文。）

为简单起见，本论证只关注 $C$ 上仅有正则奇点的情形。考虑一个单线算子。如文献 \cite{Drukker:2009tz} 与第 2.3.1 节% TODO 合稿改 \ref{subsubsec:A1-labels}
所解释，该算子的紫外标签是 $C$ 中一条非自交闭曲线 $\wp\subset C$ 的同痕类。$\langle L_\zeta(\wp)\rangle$ 的首项渐近行为将用 traffic rule 算法提取。我们选取不在反 BPS 墙上的 $(u,\zeta)$，并设 $\zeta=\vert\zeta\vert e^{i\vartheta}$，由此提取出 $\chi=\vartheta-\frac{\pi}{2}$ 对应的 WKB 三角剖分 $T_\WKB(\chi)$。（见文献 \cite{Gaiotto:2009hg} 第 6 节。）可以假定闭路径 $\wp$ 不会在 $T_\WKB(\chi)$ 的各个三角形中折返。现在考虑这一三角剖分的 traffic rule 算法。文献 \cite{Gaiotto:2009hg} 第 7.8 节证明了：对任一 WKB 三角剖分，相应的向量 $\gamma^\chi_E$（$E$ 取遍三角剖分的各边）构成一组单根基。特别地，对一切边 $E$ 都有 ${\rm Im}(e^{-i\chi}Z_{\gamma^\chi_E}(u))>0$，因此
\begin{equation}
{\rm Re}(e^{-i\vartheta}Z_{\gamma^\chi_E}(u)) > 0
\end{equation}
从而对沿半平面 $\mathbb{H}_\chi$ 趋于 $0$ 的一切 $\zeta'$，诸 $\CY_{\gamma_E^\chi}(u,\theta;\zeta')$ 的渐近行为都趋于无穷。这包括了位于其边界上的 $\zeta=e^{i\vartheta}$，所以我们仍有 $\CY_{\gamma_E^\chi}(u,\theta;\zeta)\to\infty$。由 traffic rule 现在可知，$\langle L_\zeta(\wp)\rangle$ 中的主导项简单地由
\begin{equation}
\prod_{\wp\cap E\not=\emptyset}\sqrt{\CY_{\gamma_E^\chi}}
\end{equation}
给出；于是热带荷就是 $\wp$ 所穿过的各边之荷的和：
\begin{equation}
\gamma_t(L(\wp),\check c) = \half\sum_{\wp\cap E\not=\emptyset}\gamma_E^\chi
\end{equation}

现在问：这个热带标签会怎样改变？这种改变只可能在 WKB 三角剖分发生一次 flip 时发生。\footnote{我们在此明确排除「juggle」变换，那将导致复杂得多的分析。}但文献 \cite{Gaiotto:2009hg} 的结果表明：$T_\WKB(\chi)$ 只在存在荷为 $\gamma$ 的超多重态、使得 $\chi={\rm arg}Z_\gamma(u)$ 时才改变。在当前情形 $\chi=\vartheta-\frac{\pi}{2}$，因此这意味着 $Z_\gamma/\zeta\in -i\IR_+$。也就是说，$L_\zeta(\wp)$ 的热带向量 $\gamma_t$ 只可能在反 BPS 墙上改变，而后者正是 $\check c$ 的边界。这就完成了热带理论定义中性质 1 的证明。

\insfig[0.5\textwidth]{Flip-Traffic}{利用 traffic rule 以及单根在一次 flip 下的变异，可以验证 $\gamma_t$ 跨过由一次 flip 相联系的两个腔时的热带变换规则。图中红色曲线是六个可能线算子的片段，它们不折返地穿过相应的四边形。}

现在来验证性质 2。考虑四边形中的一次 flip，以及它如何影响那些叶层 $L_1,\dots,L_6$——它们的线段与图 \ref{fig:Flip-Traffic} 中的叶层相交\zy{原文如此，疑为「与四边形相交」之笔误。}。利用文献 \cite{Gaiotto:2009hg} 的式 (7.23)--(7.24)，我们看到：当 $\chi$ 沿逆时针方向经过某个临界相位时，诸单根按一次变异而改变：
\begin{equation}
\begin{split}
\gamma_{E^+}^+ & = -\gamma_{E^-}^- \\
\gamma_{E_i^+}^+ & = \gamma_{E_i^-}^- + \langle\gamma_{E_i^-}^-,\gamma_{E^-}^-\rangle_+ \gamma_{E^-}^-
\end{split}
\end{equation}
由此容易检验热带变换规则 \eqref{eq:Tropical-Transformation}。

\insfig[0.5\textwidth]{tropic-onto-A}{乘上具有纯味荷的适当线算子后，上述线算子的热带标签为 $\gamma_E$。}

\insfig[0.5\textwidth]{tropic-onto-B}{乘上具有纯味荷的适当线算子后，上述线算子的热带标签为 $-\gamma_E$。}

为确立性质 3，只需给出热带标签分别等于 $\pm\gamma^\chi_E$ 的单线算子。于是，借助 \eqref{eq:tropic-mult-line}，这类线算子的乘积将给出任意想要的向量。现在参照图 \ref{fig:tropic-onto-A}，我们看到路径 $\wp$ 的热带标签为
\begin{equation}
\gamma_t(L(\wp)) = \left(\half\sum_{v_2\in E_i}\gamma_{E_i} + \half\gamma_E\right) + \left(\half\sum_{v_4\in E_i}\gamma_{E_i} + \half\gamma_E\right)
\end{equation}
其次注意：绕顶点 $v$、权重为 $-1$ 的小圆 $C_v$ 所构成的叶层满足 $L(C_v)=\mu_v^{-1}$，因而热带标签为 $-\half\sum_{v\in E_i}\gamma_{E_i}$。因此 $L(C_{v_1})L(C_{v_2})L(\wp)$ 的热带标签为 $\gamma_E$。类似地，利用图 \ref{fig:tropic-onto-B} 中的路径 $\wp'$，我们发现 $L(C_{v_1})L(C_{v_2})L(\wp')$ 的热带标签为 $-\gamma_E$。

最后，为确立性质 4，设给定一个热带向量 $\gamma_t$。由于诸 $\gamma_{E_i}$ 是单根，它具有到 $\gamma_{E_i}$ 之和的唯一分解，于是我们知道它穿过了哪些边。现在考虑 WKB 三角剖分的对偶胞腔分解。对热带向量中出现的每一条边 $E_i$（按重数计），我们在对偶胞腔分解中配上一条边。这些边以唯一的方式拼合成一个叶层，从而完成「$A_1$ 理论（仅有正则奇点）是热带的」这一证明。我们预期，一个类似但也许更繁复的论证将覆盖带不规则奇点的情形。

本节末尾我们作几点评注：

\begin{enumerate}
\item 变换 \eqref{eq:Tropical-Transformation} 是 Kontsevich--Soibelman 变换 $\CK_{\gamma_0}^{\Omega(\gamma_0)}$ 的「热带」版本。我们可以把它记作 $\CK_{t,\gamma_0}^{\Omega(\gamma_0)}$。存在一个热带穿墙公式，其中 $\CK$ 处处被换成 $\CK_t$。注意，上面我们只考虑了与 BPS 超多重态相联系的 flip；没有考虑与矢量多重态相联系的、更为困难的「juggle」变换（见文献 \cite{Gaiotto:2009hg} 第 6.6.3 节）。

\item 这里与斯托克斯理论有密切的类比。鉴于文献 \cite{Gaiotto:2008cd} 的 $\zeta$-微分方程，这并非偶然。

\item 在 $A_1$ 理论中，$L(\wp)$ 的热带向量 $\gamma_t$ 有一个非常漂亮的物理解释：它是 Seiberg--Witten 曲线 $\Sigma$ 中主导渐近行为的那条 WKB 路径的同调类。
\end{enumerate}

\subsection{例：$N=3$ Argyres--Douglas 理论的热带标签}

在式 (10.6)% TODO 合稿改 \eqref{eq:ADN3-gen-lam}
中，我们写下了 $N=3$ AD 理论中一般的单叶层（simple lamination）。这会把单叶层的集合分解成五个「象限」，由 $m,n\ge 0$ 给出坐标。事实上，把这五个象限粘成单独一份 $\IZ^2$ 是很自然的，如下所述。给定任一单线算子 $L$ 与任一带有坐标 $X$、$Y$ 的三角剖分，我们把 $\inprod{L}'$ 按 $X$ 与 $Y$ 展开。结果总具有形式
\begin{equation}
\inprod{L}' = X^a Y^b P(1/X,1/Y)
\end{equation}
其中 $P=1+\cdots$ 是某个多项式。数对 $(a,b)$ 就是算子 $L$ 的热带向量（相对于 $\Gamma$ 的基 $\gamma_X$、$\gamma_Y$）。

例如，用我们上面引入的三角剖分，五个象限中的算子 $L_i^m L_{i+1}^n$ 分别映射到 $(n+m,n)$、$(n,m+n)$、$(-n,m)$、$(-m,-n)$、$(n,-m)$；注意这五个区域合起来填满 $\IZ^2$，所以每种可能的红外荷都有相应的单线算子。这其实只是来自这样一个事实：五个典范的 $L_i$ 按循环次序排列时，其首项依次为 $X$、$XY$、$Y$、$1/X$、$1/Y$，而它们在单项式平面上也是按循环次序排列的。

对每个三角剖分，我们就这样得到所有单线算子之空间的一个自然坐标化，把它等同于 $\IZ^2$。对三角剖分作 flip 则给出这些坐标的一个简单的分段线性变换。这一结构在文献 \cite{qd-pentagon} 中有非常漂亮的描述。
